% Introducció.
%
% Estructura recomanada:
%   1. Motivació: per què aquest problema és interessant o necessari.
%   2. Objectius: objectiu general + objectius específics (llista numerada).
%   3. Abast: què queda dins i què queda fora del treball.
%   4. Estructura del document: un paràgraf breu resumint què conté cada capítol.


Aquest Treball de Fi de Grau (TFG) planteja una anàlisi comparativa d'orquestració de contenidors en entorns \textit{Edge}: s'hi contrasta una solució d'execució nativa basada en Systemd amb un orquestrador estàndard com Kubernetes (en la seva variant lleugera, K3s).

\section{Motivació}
Els dispositius perimetrals acostumen a operar amb restriccions de maquinari notables. Aquesta limitació topa amb el desplegament d'orquestradors moderns com Kubernetes, els quals requereixen l'execució permanent de processos secundaris \textemdash{}com ara el \textit{kubelet} o els components de xarxa\textemdash{} per garantir el seu funcionament. Aquest consum estructural actua com un llast de memòria i CPU, reduint el marge disponible per a càrregues de treball exigents, com per exemple models d'Intel·ligència Artificial. Per tant, resulta clau explorar alternatives d'execució directa a nivell de sistema operatiu per determinar si és possible recuperar aquesta capacitat de còmput sense comprometre la fiabilitat del sistema.

\section{Objectius}
L'objectiu principal d'aquest projecte és avaluar el sobrecost de rendiment que comporta l'ús de Kubernetes (K3s) en maquinari de capacitats reduïdes quan s'hi executen tasques d'Intel·ligència Artificial (concretament, l'entrenament d'un model MNIST amb PyTorch). Aquest escenari es compara amb una arquitectura nativa gestionada mitjançant Systemd, Ansible i Podman.

Per assolir aquesta premissa, es fixen els objectius específics següents:
\begin{itemize}
    \item Quantificar l'impacte sobre la infraestructura mitjançant l'anàlisi de paràmetres base, tals com el temps de desplegament (\texttt{T\_deploy}), el consum de CPU en repòs i l'ocupació de memòria RAM base.
    \item Avaluar el rendiment de xarxa sotmetent ambdues arquitectures a proves d'estrès sota quatre protocols diferents: HTTP-JSON, gRPC-Protobuf, ZeroMQ-Protobuf i ZeroMQ-MessagePack.
    \item Analitzar la penalització de la xarxa virtual a través de mètriques com el temps total de transmissió, el \textit{throughput}, la taxa d'èxit (\%), la latència RTT, el temps de processament (\texttt{T\_Proc}), els pics de recursos i la mida de la càrrega útil (\textit{payload}).
    \item Determinar la capacitat d'autorecuperació davant d'errors (MTTR) mitjançant la terminació abrupta de processos, contrastant la resiliència del model natiu com la de l'orquestrador.
    \item Automatitzar completament tant els procediments de desplegament com l'extracció de dades mitjançant \textit{scripts} dedicats, assegurant així la total reproductibilitat de l'experiment.
\end{itemize}

\section{Abast}
L'estudi se centra en l'avaluació de nodes \textit{Edge} aïllats, contraposant la versió lleugera de Kubernetes (K3s) amb una integració de Podman i Systemd basada en \textit{Quadlets}, on cada \textit{worker} executa càrregues d'entrenament del model MNIST amb PyTorch. Aquesta configuració s'acompanya de l'execució sistemàtica de proves de rendiment sobre els protocols de comunicació indicats.

L'avaluació quantitativa s'estructura en dos eixos principals: d'una banda, les mètriques d'infraestructura destinades a mesurar l'impacte base en l'arrencada i el consum de recursos (incloent-hi l'MTTR); de l'altra, les mètriques de xarxa orientades a caracteritzar l'eficiència comunicativa sota condicions de càrrega.

Queden excloses de l'abast les distribucions tradicionals de Kubernetes pròpies d'entorns \textit{Cloud}, el disseny d'arquitectures d'alta disponibilitat (HA) per al pla de control, i la utilització d'altres motors de contenidors basats en dimoni, com ara Docker.

\section{Estructura del document}
El document s'organitza en set capítols que descriuen de manera progressiva el desenvolupament i l'avaluació del projecte:

\begin{description}
    \item[Capítol 1 (Introducció):] Introdueix el context de la recerca, la motivació, els objectius establerts i l'abast tècnic.
    
    \item[Capítol 2 (Marc Teòric):] Recull els fonaments tecnològics sobre la gestió de contenidors, els protocols de xarxa avaluats i les eines emprades (Systemd, Ansible, Kubernetes).
    
    \item[Capítol 3 (Estat de l'Art):] Examina les solucions actuals de la indústria davant de la gestió de càrregues de treball en dispositius \textit{Edge} amb recursos limitats.
    
    \item[Capítol 4 (Metodologia):] Detalla el disseny experimental, l'entorn de proves configurat i la definició formal de les mètriques monitorades.
    
    \item[Capítol 5 (Implementació):] Descriu el procés pràctic de desplegament, la integració mitjançant \textit{Quadlets} i l'automatització dels \textit{scripts} de mesura.
    
    \item[Capítol 6 (Resultats):] Presenta l'anàlisi de les dades obtingudes, contrastant el comportament d'ambdues arquitectures pel que fa al consum de recursos, rendiment de xarxa i resiliència.
    
    \item[Capítol 7 (Conclusions):] Sintetitza el compliment dels objectius, avalua els compromisos tècnics (\textit{trade-offs}) observats i proposa possibles línies de continuïtat per a futurs treballs.
\end{description}